\chapter{Funções características e o teorema central do
limite}\label{ch:b3:clt}

A lei dos grandes números diz que as médias convergem; o teorema central
do limite diz \emph{como elas flutuam}: o erro, ampliado por $\sqrt n$, é
assintoticamente \omterm{def:b3:clt:gaussianvector}{gaussiano} --- qualquer que seja a lei de partida. Essa
universalidade é o fato mais profundo da probabilidade elementar, e sua
demonstração natural é de análise de Fourier: a \emph{\omterm{def:b3:clt:cf}{função
característica}} (a transformada de Fourier de uma lei) converte somas
\omterm{def:b3:probability:independence}{independentes} em produtos, e o maquinário do~\cref{ch:b3:fouriertransform} ---
injetividade, pontos fixos \omterm{def:b3:clt:gaussianvector}{gaussianos} --- converte a convergência
pontual desses produtos em convergência de leis (teorema de Lévy,
demonstrado por inteiro). O capítulo termina com os \omterm{def:b3:clt:gaussianvector}{vetores gaussianos} e
a dedução honesta dos intervalos de confiança usados em toda a
estatística; o problema de fim de semana dá a segunda demonstração do
teorema central do limite, a de Lindeberg, com taxa de erro explícita.

\section{Funções características}

\begin{definition}\label{def:b3:clt:cf}
A \emph{função característica}\index{função característica} de uma
\omterm{def:b3:probability:space}{variável aleatória} real $X$ é
\[
\varphi_X(\xi) = \E\bigl[\eu^{\iu\xi X}\bigr]
= \int_\R \eu^{\iu\xi x}\,\dd\P_X(x)
\qquad (\xi \in \R)
\]
(o teorema de transferência a calcula a partir da lei; para uma
\omterm{ex:b3:lebesgue:gamma}{densidade} $f$, $\varphi_X(\xi) = \hat f(-\xi)$ na convenção do~\cref{ch:b3:fouriertransform}).
\end{definition}

\begin{proposition}\label{prop:b3:clt:cfbasics}
(a) $\varphi_X(0) = 1$, $\abs{\varphi_X} \leq 1$, e $\varphi_X$ é uniformemente \omterm{def:b3:topology:continuity}{contínua}; $\varphi_{aX + b}(\xi) = \eu^{\iu b\xi}\varphi_X(a\xi)$.
(b) Se $X, Y$ são \emph{\omterm{def:b3:probability:independence}{independentes}}: $\varphi_{X+Y} = \varphi_X\,\varphi_Y$.
(c) Se $\E\abs X^k < \infty$, então $\varphi_X \in \mathcal
C^k$ com $\varphi_X^{(j)}(0) = \iu^j\,\E[X^j]$ para $j \leq
k$; em particular, para $X \in L^2$
centrada de variância $\sigma^2$:
\[
\varphi_X(\xi) = 1 - \frac{\sigma^2\xi^2}{2} +
o(\xi^2) \qquad (\xi \to 0).
\]
(d) \omterm{def:b3:clt:gaussianvector}{Gaussiana}: $X \sim \mathcal N(m, \sigma^2)$ tem $\varphi_X(\xi) = \eu^{\iu m\xi - \sigma^2\xi^2/2}$.
\end{proposition}

\begin{proof}
(a) As cotas são imediatas; \omterm{def:b3:topology:continuity}{continuidade}: $\abs{\varphi(\xi + h)
- \varphi(\xi)} \leq \E\abs{\eu^{\iu hX} - 1} \to 0$ quando $h
\to 0$, por
convergência dominada, uniformemente em $\xi$. A regra afim é uma
substituição. (b) $\eu^{\iu\xi(X+Y)} =
\eu^{\iu\xi X}\eu^{\iu\xi Y}$, e as \omterm{def:b3:probability:space}{esperanças} de produtos de variáveis
\omterm{def:b3:probability:independence}{independentes} se fatorizam (\cref{thm:b3:probability:independence}, aplicado às partes real e
imaginária). (c) Derivação sob a \omterm{def:b3:probability:space}{esperança}, dominada por $\E\abs X^j$
(\cref{thm:b3:lebesgue:paramdiff}); a expansão de Taylor em $0$ é então
Taylor--Young para a função $\varphi$ de classe $\mathcal C^2$. (d) Para $\mathcal N(0,1)$: a
transformada da \omterm{def:b3:clt:gaussianvector}{gaussiana} (\cref{ex:b3:fouriertransform:gaussian} com $a = \frac12$) dá $\int\eu^{\iu\xi x}\frac{\eu^{-x^2/2}}{\sqrt{2\pi}}\dd
x = \eu^{-\xi^2/2}$; o caso geral,
pela regra afim.
\end{proof}

\begin{theorem}[Injetividade]\label{thm:b3:clt:injectivity}
Se $\varphi_X = \varphi_Y$, então $X$ e $Y$ têm a mesma lei. Mais precisamente, para
$N \sim \mathcal N(0,1)$ \omterm{def:b3:probability:independence}{independente} de $X$ e $\varepsilon > 0$, a variável suavizada $X +
\varepsilon N$ tem a
\omterm{ex:b3:lebesgue:gamma}{densidade}
\[
p_\varepsilon(x) = \frac1{2\pi}\int_\R
\varphi_X(-\xi)\,\eu^{-\varepsilon^2\xi^2/2}\,
\eu^{\iu\xi x}\,\dd\xi ,
\]
determinada apenas por $\varphi_X$; fazer $\varepsilon \to 0$ recupera a lei de
$X$.\index{função característica, injetividade}
\end{theorem}

\begin{proof}
$X + \varepsilon N$ tem \omterm{ex:b3:lebesgue:gamma}{densidade} $p_\varepsilon(x) =
\E\bigl[g_\varepsilon(x - X)\bigr]$, em que $g_\varepsilon$ é a \omterm{ex:b3:lebesgue:gamma}{densidade} $\mathcal N(0, \varepsilon^2)$: com efeito, para
$B$ boreliano, a \omterm{def:b3:probability:independence}{independência} e Tonelli dão $\P(X + \varepsilon N \in
B) = \int\!\!\int\mathbf 1_B(x + \varepsilon
n)g_1(n)\,\dd n\,\dd\P_X(x) = \int_B\E[g_\varepsilon(t -
X)]\dd t$ (substitua e, depois,
Tonelli de novo). Escrevendo $g_\varepsilon$ pela inversão de Fourier de sua
transformada (\cref{exo:b3:fouriertransform:4}, reescalada): $g_\varepsilon(u) = \frac1{2\pi}\int
\eu^{-\varepsilon^2\xi^2/2}\eu^{\iu\xi u}\dd\xi$, e Fubini (tudo dominado
pelo fator \omterm{def:b3:clt:gaussianvector}{gaussiano}):
\[
p_\varepsilon(x) = \frac1{2\pi}\int_\R
\varphi_X(-\xi)\,\eu^{-\varepsilon^2\xi^2/2}\,
\eu^{\iu\xi x}\,\dd\xi ,
\]
um funcional apenas de $\varphi_X$. Se $\varphi_X =
\varphi_Y$: $X + \varepsilon N$ e $Y + \varepsilon N$ têm leis iguais para
todo $\varepsilon$; para $f$ \omterm{def:b3:topology:continuity}{contínua} limitada, $\E f(X + \varepsilon N) \to \E f(X)$ quando $\varepsilon
\to 0$ (convergência
dominada, $X + \varepsilon N \to X$ pontualmente no espaço produto), de modo que $\E f(X) = \E f(Y)$ para
toda tal $f$ --- e isso determina a lei: para cada $t$, comprima
$\mathbf 1_{\intoc{-\infty}t}$ entre as rampas \omterm{def:b3:topology:continuity}{contínuas} limitadas $f_k^\pm$ (iguais a $1$ em $\intoc{-\infty}{t \mp \frac1k}$, a
$0$ além de $t \pm
\frac1k$, e afins entre esses trechos); passando ao limite em
$\E
f_k^-(X) \leq F_X(t) \leq \E f_k^+(X)$ obtém-se $F_X(t) =
F_Y(t)$ em todo $t$ em que ambas são \omterm{def:b3:topology:continuity}{contínuas}, logo em
toda parte, por \omterm{def:b3:topology:continuity}{continuidade} à direita e pela \omterm{ex:b3:lebesgue:gamma}{densidade} dos pontos de
\omterm{def:b3:topology:continuity}{continuidade} comuns (as duas $F$ têm um número enumerável de saltos);
funções de distribuição iguais forçam leis iguais (\cref{exo:b3:measure:3}, que
repousa sobre o~\cref{thm:b3:measure:uniqueness}).
\end{proof}

\section{Convergência em distribuição}

\begin{definition}\label{def:b3:clt:cid}
$X_n$ \emph{converge em distribuição} (ou em lei) para $X$, o que
se escreve $X_n \Rightarrow X$, se\index{convergência em distribuição}
\[
\E\bigl[f(X_n)\bigr] \longrightarrow \E\bigl[f(X)\bigr]
\qquad\text{para toda contínua limitada } f\colon\R\to\R .
\]
Equivalentemente (\cref{exo:b3:clt:4}): $F_{X_n}(t) \to F_X(t)$ em todo ponto de \omterm{def:b3:topology:continuity}{continuidade}
$t$ de $F_X$. Não é preciso que as $X_n$ vivam em um \omterm{def:b3:probability:space}{espaço de
probabilidade} comum: só as leis importam.
\end{definition}

\begin{theorem}[Teorema de seleção de Helly]
\label{thm:b3:clt:helly}
Toda sequência $(F_n)$ de funções de distribuição tem uma
subsequência que converge pontualmente, em todo ponto de \omterm{def:b3:topology:continuity}{continuidade} do
limite, para uma função $G \colon
\R \to \intcc01$ não decrescente e \omterm{def:b3:topology:continuity}{contínua} à direita ---
possivelmente com $G(+\infty) - G(-\infty) <
1$ (a massa pode escapar para o infinito).\index{teorema
de Helly}
\end{theorem}

\begin{proof}
Uma extração diagonal dá $F_{n_k}(q) \to \ell(q)$ para todo racional $q$ (valores no
\omterm{def:b3:topology:compact}{compacto} $\intcc01$). Defina $G(t) = \inf\{\ell(q) : q \in \Q, q > t\}$: não decrescente; \omterm{def:b3:topology:continuity}{contínua} à direita (um
ínfimo sobre \omterm{def:b3:topology:topology}{vizinhanças} racionais que encolhem pela direita). Em um
ponto de \omterm{def:b3:topology:continuity}{continuidade} $t$ de $G$: para racionais $q_1 < t < q_2$,
\[
\ell(q_1) \leq \liminf F_{n_k}(t) \leq \limsup F_{n_k}(t)
\leq \ell(q_2),
\]
pela monotonicidade de cada $F_{n_k}$. Da definição de $G$ como ínfimo e
da monotonicidade de $\ell$ nos racionais: $G(s) \leq \ell(q) \leq G(q)$ sempre que $s < q$.
Tomando $s < q_1 < t$ obtém-se $\ell(q_1) \geq G(s)$ e $\ell(q_2) \leq G(q_2)$; fazendo $s \uparrow t$ e $q_2
\downarrow t$, a \omterm{def:b3:topology:continuity}{continuidade}
de $G$ em $t$ comprime tanto o $\liminf$ quanto o $\limsup$ até $G(t)$.
\end{proof}

\begin{lemma}[Tensão a partir da função característica]
\label{lem:b3:clt:tightness}
Para uma \omterm{def:b3:probability:space}{variável aleatória} $X$ qualquer e $u > 0$:
\[
\P\Bigl(\abs X \geq \frac2u\Bigr) \;\leq\;
\frac1u\int_{-u}^{u}\bigl(1 -
\operatorname{Re}\varphi_X(\xi)\bigr)\,\dd\xi .
\]
\end{lemma}

\begin{proof}
Por Tonelli--Fubini (integrando limitado, região finita em $\xi$):
\[
\frac1u\int_{-u}^u\bigl(1 -
\operatorname{Re}\varphi_X(\xi)\bigr)\dd\xi
= \E\Bigl[\frac1u\int_{-u}^u(1 - \cos(\xi X))\,\dd\xi\Bigr]
= 2\,\E\Bigl[1 - \frac{\sin(uX)}{uX}\Bigr]
\]
(interprete o colchete como seu limite $0$ em $X = 0$). O
integrando é não negativo ($\abs{\sin t} \leq \abs t$) e, para $\abs{uX} \geq 2$: $1 - \frac{\sin(uX)}{uX} \geq 1 -
\frac1{\abs{uX}} \geq \frac12$. Manter apenas o
evento $\{\abs{uX} \geq 2\}$ dentro da \omterm{def:b3:probability:space}{esperança} deixa, portanto, ao menos $2 \cdot \frac12\,\P(\abs X \geq \frac2u)$, que é
o que se afirmava.
\end{proof}

\begin{theorem}[Teorema da continuidade de Lévy]
\label{thm:b3:clt:levy}
Sejam $(X_n)$ \omterm{def:b3:probability:space}{variáveis aleatórias} cujas \omterm{def:b3:clt:cf}{funções características}
convergem pontualmente: $\varphi_{X_n}(\xi) \to
\varphi(\xi)$ para todo $\xi$, em que $\varphi =
\varphi_X$ é a
\omterm{def:b3:clt:cf}{função característica} de alguma \omterm{def:b3:probability:space}{variável aleatória} $X$. Então
$X_n \Rightarrow X$.\index{teorema da continuidade de Lévy}
\end{theorem}

\begin{proof}
\emph{Tensão.} Fixe $\varepsilon > 0$. Como $\varphi$ é \omterm{def:b3:topology:continuity}{contínua} em $0$ com $\varphi(0) = 1$,
escolha $u > 0$ com $\frac1u\int_{-u}^u(1 - \operatorname{Re}\varphi) <
\varepsilon$; por convergência dominada (o integrando é
limitado por $2$ no $[-u,u]$ fixado), a mesma integral para $\varphi_{X_n}$
é $< 2\varepsilon$ para $n$ grande: o~\cref{lem:b3:clt:tightness} dá $\P(\abs{X_n} \geq \frac2u)
\leq 2\varepsilon$ para $n$
grande, e aumentar a constante cuida das finitas outras: as leis são
\emph{tensas} --- nenhuma massa escapa.

\emph{Subsequências.} Seja $(F_{n_k})$ uma subsequência qualquer; por Helly
(\cref{thm:b3:clt:helly}), extraia $F_{n_{k_j}} \to G$ nos pontos de \omterm{def:b3:topology:continuity}{continuidade}. A tensão
força $G(-\infty) = 0$, $G(+\infty) = 1$ ($G(\frac2u) - G(-\frac2u) \geq 1 -
2\varepsilon$ nos pontos de \omterm{def:b3:topology:continuity}{continuidade}): $G$ é uma
função de distribuição genuína, de alguma \omterm{def:b3:probability:space}{variável aleatória} $Y$.
Então $X_{n_{k_j}} \Rightarrow Y$ (\cref{exo:b3:clt:4}, convergência em distribuição a partir das
$F$), de modo que $\varphi_{X_{n_{k_j}}} \to \varphi_Y$ \emph{pontualmente} ($x
\mapsto \eu^{\iu\xi x}$ é \omterm{def:b3:topology:continuity}{contínua} e
limitada, nas partes real e imaginária separadamente); comparando com a
hipótese: $\varphi_Y = \varphi = \varphi_X$, e a injetividade (\cref{thm:b3:clt:injectivity}) dá $Y \sim X$, isto é,
$G =
F_X$.

\emph{Conclusão.} Toda subsequência de $(F_n)$ tem uma
subsubsequência que converge à \emph{mesma} $F_X$ (em seus pontos de
\omterm{def:b3:topology:continuity}{continuidade}); logo $F_n(t) \to F_X(t)$ em todo ponto de \omterm{def:b3:topology:continuity}{continuidade} $t$ (uma
sequência real cujas subsequências todas têm subsubsequências com o
mesmo limite converge): $X_n \Rightarrow X$.
\end{proof}

\section{O teorema central do limite}

\begin{theorem}[Teorema central do limite]\label{thm:b3:clt:clt}
Sejam $(X_n)$ i.i.d.\ com $\E X_1 = m$ e $\V(X_1) =
\sigma^2 \in \intoo0\infty$. Então\index{teorema central
do limite}
\[
\frac{S_n - nm}{\sigma\sqrt n} \;\Longrightarrow\; \mathcal
N(0, 1) :
\qquad
\P\Bigl(a \leq \frac{S_n - nm}{\sigma\sqrt n} \leq
b\Bigr) \longrightarrow
\frac{1}{\sqrt{2\pi}}\int_a^b\eu^{-x^2/2}\,\dd x
\]
para todos $a < b$.
\end{theorem}

\begin{proof}
Centre e normalize: $Z_i = \frac{X_i - m}{\sigma}$ (i.i.d., de média $0$ e variância $1$)
e $T_n =
\frac1{\sqrt n}\sum_{i\leq n}Z_i$. Pela \omterm{def:b3:probability:independence}{independência} e pela regra afim (\cref{prop:b3:clt:cfbasics}):
\[
\varphi_{T_n}(\xi) =
\varphi_{Z}\Bigl(\frac{\xi}{\sqrt n}\Bigr)^{n},
\qquad
\varphi_Z(\eta) = 1 - \frac{\eta^2}2 + \eta^2\rho(\eta),\quad
\rho(\eta)\to0 .
\]
Fixe $\xi$ e ponha $a_n = \varphi_Z(\xi/\sqrt n)$, $b_n = 1 -
\frac{\xi^2}{2n}$: ambos têm módulo $\leq 1$ para
$n$ grande ($\abs{b_n} \leq 1$ assim que $\xi^2 \leq 4n$; $\abs{a_n} \leq 1$ sempre). A desigualdade
elementar $\abs{a^n - b^n} \leq
n\abs{a - b}$ para $\abs a, \abs b \leq 1$ (telescopando $a^n - b^n = \sum a^k(a - b)b^{n-1-k}$) dá
\[
\Bigl|\varphi_{T_n}(\xi) - \Bigl(1 -
\frac{\xi^2}{2n}\Bigr)^{n}\Bigr|
\leq n\,\Bigl|\varphi_Z\Bigl(\frac\xi{\sqrt n}\Bigr) - 1 +
\frac{\xi^2}{2n}\Bigr|
= \xi^2\,\Bigl|\rho\Bigl(\frac{\xi}{\sqrt n}\Bigr)\Bigr|
\longrightarrow 0,
\]
ao passo que $\bigl(1 - \frac{\xi^2}{2n}\bigr)^n \to
\eu^{-\xi^2/2}$ (logaritmo real). Logo $\varphi_{T_n}(\xi) \to
\eu^{-\xi^2/2} = \varphi_{\mathcal N(0,1)}(\xi)$ (\cref{prop:b3:clt:cfbasics}(d))
para todo $\xi$: Lévy (\cref{thm:b3:clt:levy}) conclui $T_n \Rightarrow \mathcal
N(0,1)$. As
probabilidades de intervalos decorrem disso, pois $F_{\mathcal
N}$ é \omterm{def:b3:topology:continuity}{contínua} em
toda parte.
\end{proof}

\begin{example}[Intervalos de confiança, deduzidos honestamente]
\label{ex:b3:clt:confidence}
Consulte $n$ eleitores \omterm{def:b3:probability:independence}{independentes}; $\hat p_n = S_n/n$ estima o $p$
verdadeiro, com $\sigma^2 = p(1-p) \leq \frac14$. O teorema central do limite dá, para $n$
grande,
\[
\P\Bigl(\abs{\hat p_n - p} \leq
\frac{z}{2\sqrt n}\Bigr)
\;\geq\; \P\Bigl(\Bigl|\frac{S_n - np}{\sigma\sqrt n}\Bigr|
\leq z\Bigr)
\longrightarrow \Phi(z) - \Phi(-z),
\]
em que $\Phi$ é a função de distribuição \omterm{def:b3:clt:gaussianvector}{gaussiana} padrão. Com
$z = 1.96$: confiança assintótica de $95\%$, e uma margem $\frac{1.96}{2\sqrt n} \leq 3\%$ exige
$n \geq
\bigl(\frac{1.96}{0.06}\bigr)^2 \approx 1068$ --- o número por trás de todo ``$\pm3$ pontos, $95\%$'' que
se lê por aí; compare com os $5556$ de Chebyshev
(\cref{exo:b3:probability:7}). O $\sqrt n$ é universal: para reduzir o erro à metade,
quadruplique a amostra --- a mesma lei que fixa o custo do \omterm{ex:b3:probability:sllnapps}{Monte Carlo}
(\cref{exo:b3:clt:7}).
\end{example}

\section{Vetores gaussianos}

\begin{definition}\label{def:b3:clt:gaussianvector}
Um vetor aleatório $X = (X_1, \dots, X_d)$ é \emph{gaussiano}\index{vetor gaussiano} se
toda combinação linear $\langle t, X\rangle = \sum t_iX_i$ é uma variável gaussiana real (possivelmente
degenerada). Sua lei é determinada pelo vetor de médias $m = (\E X_i)$ e pela
\emph{matriz de covariância} $\Sigma = \bigl(\operatorname{Cov}
(X_i, X_j)\bigr)$: com efeito, a \omterm{def:b3:clt:cf}{função característica}
do vetor, $\varphi_X(t) = \E\eu^{\iu\langle t, X\rangle}$, é o valor em $1$ da \omterm{def:b3:clt:cf}{função característica} de $\langle t, X\rangle$:
\[
\varphi_X(t) = \exp\Bigl(\iu\langle t, m\rangle -
\tfrac12\,t^{\mathsf T}\Sigma\,t\Bigr),
\]
e as \omterm{def:b3:clt:cf}{funções características} em dimensão $d$ são injetoras (a mesma
demonstração por suavização do~\cref{thm:b3:clt:injectivity}, com gaussianas coordenada
a coordenada).
\end{definition}

\begin{theorem}\label{thm:b3:clt:gaussianvector}
Seja $X$ um \omterm{def:b3:clt:gaussianvector}{vetor gaussiano}.
\begin{enumerate}
  \item Toda imagem afim $AX + b$ é um \omterm{def:b3:clt:gaussianvector}{vetor gaussiano}.
  \item As componentes $X_i$ são \emph{\omterm{def:b3:probability:independence}{independentes}} se, e somente
        se, $\Sigma$ é diagonal: para variáveis conjuntamente
        \omterm{def:b3:clt:gaussianvector}{gaussianas}, não correlacionadas $=$ \omterm{def:b3:probability:independence}{independentes}.
  \item Se $\Sigma$ é inversível, $X$ tem \omterm{ex:b3:lebesgue:gamma}{densidade} $\frac{1}{(2\pi)^{d/2}\sqrt{\det\Sigma}}
        \exp\bigl(-\frac12(x - m)^{\mathsf T}\Sigma^{-1}(x -
        m)\bigr)$.
\end{enumerate}
\end{theorem}

\begin{proof}
(1) As combinações lineares das componentes de $AX + b$ são funções
afins de combinações lineares de $X$: \omterm{def:b3:clt:gaussianvector}{gaussianas} (uma imagem afim de
uma variável \omterm{def:b3:clt:gaussianvector}{gaussiana} é \omterm{def:b3:clt:gaussianvector}{gaussiana}).
(2) Se $\Sigma$ é diagonal, a \omterm{def:b3:clt:cf}{função característica} se fatoriza:
$\varphi_X(t) = \prod_i\exp(\iu t_im_i -
\frac12\Sigma_{ii}t_i^2) = \prod\varphi_{X_i}(t_i)$, que é a \omterm{def:b3:clt:cf}{função característica} da lei produto (\cref{thm:b3:probability:independence} lido
através da injetividade em dimensão $d$): as componentes são
\omterm{def:b3:probability:independence}{independentes}. A recíproca é o anulamento das covariâncias de variáveis
$L^2$ \omterm{def:b3:probability:independence}{independentes}.
(3) Diagonalize $\Sigma = P D P^{\mathsf T}$ ($P$ ortogonal, $D > 0$ diagonal ---
\cref{exo:b3:submanifolds:8}); o vetor $Y = P^{\mathsf T}(X - m)$ é \omterm{def:b3:clt:gaussianvector}{gaussiano} de covariância $D$: por
(2), suas componentes são $\mathcal N(0, d_i)$ \omterm{def:b3:probability:independence}{independentes}, de modo que $Y$ tem a
\omterm{ex:b3:lebesgue:gamma}{densidade} produto; empurre-a para a frente pela aplicação $x = m + PY$, que
preserva volume (\cref{thm:b3:product:linearchange}, $\abs{\det P} = 1$), e reescreva o expoente de
maneira invariante.
\end{proof}

\begin{theorem}[Teorema central do limite multidimensional]
\label{thm:b3:clt:multiclt}
Sejam $(X_n)$ \emph{vetores} aleatórios i.i.d.\ de $\R^d$, de
quadrado \omterm{def:b3:lebesgue:l1}{integrável}, com média $m$ e matriz de covariância
$\Sigma$. Então $\frac{S_n - nm}{\sqrt n}$ converge em distribuição ao \omterm{def:b3:clt:gaussianvector}{vetor gaussiano}
$\mathcal N(0,
\Sigma)$.\index{teorema central do limite!multidimensional}
\admitted
\end{theorem}

\begin{remark}\label{rem:b3:clt:multiclt}
Quase tudo já está em nossas mãos. Para cada direção $t \in \R^d$, a variável
real $\langle t,
\frac{S_n - nm}{\sqrt n}\rangle$ é uma soma normalizada de variáveis reais i.i.d.\ de variância
$t^{\mathsf T}\Sigma t$, de modo que o cálculo do~\cref{thm:b3:clt:clt} dá a convergência pontual
das \omterm{def:b3:clt:cf}{funções características} em dimensão $d$ para $\eu^{-t^{\mathsf T}\Sigma t/2}$, a \omterm{def:b3:clt:cf}{função
característica} de $\mathcal N(0, \Sigma)$ (\cref{def:b3:clt:gaussianvector}). O que não redemonstramos é o
teorema da \omterm{def:b3:topology:continuity}{continuidade} de Lévy \emph{em $\R^d$}: a seleção de Helly
e a estimativa de tensão se generalizam rotineiramente (coordenada a
coordenada), e essa redução de Cramér--Wold é feita honestamente em
qualquer curso de pós-graduação em probabilidade; nada além dos métodos
deste capítulo é necessário.
\end{remark}

\begin{method}\label{met:b3:clt:method}
Para identificar uma lei-limite: calcule \omterm{def:b3:clt:cf}{funções características}, tome o
limite pontual, reconheça-o (\omterm{def:b3:clt:gaussianvector}{gaussiana} $\eu^{-\sigma^2\xi^2/2}$, Poisson $\eu^{\lambda(\eu^{\iu\xi}-1)}$, exponencial
$\frac{\lambda}
{\lambda - \iu\xi}$, \dots) e invoque Lévy. O ritual de três passos (\omterm{def:b3:probability:independence}{independência}
$\to$ produto; Taylor em $0$ $\to$ limite exponencial; Lévy
$\to$ convergência em lei) demonstra o teorema central do limite, a
lei dos eventos raros de Poisson (\cref{exo:b3:clt:5}) e todo teorema-limite
clássico deste curso. Para enunciados q.c., volte à caixa de ferramentas
do~\cref{ch:b3:probability}: os dois capítulos respondem a perguntas diferentes
sobre o mesmo $S_n$.
\end{method}

\section{Exercícios}

\begin{exercise}[$\star$]\label{exo:b3:clt:1}
Calcule as \omterm{def:b3:clt:cf}{funções características}: uniforme em $\intcc{-1}1$; exponencial
$\mathcal E(\lambda)$; Poisson $\mathcal P(\lambda)$; binomial $\mathcal B(n, p)$. Deduza pelo~\cref{thm:b3:clt:injectivity} que a
soma de variáveis de Poisson \omterm{def:b3:probability:independence}{independentes} ($\lambda, \mu$) é Poisson $(\lambda +
\mu)$.
\end{exercise}

\begin{exercise}[$\star\star$]\label{exo:b3:clt:2}
(a) Mostre que $\varphi_X$ toma valores reais se, e somente se, $X$ e
$-X$ têm a mesma lei (uma variável \emph{simétrica}).
(b) Suponha $\abs{\varphi_X(\xi_0)} = 1$ para algum $\xi_0 \neq
0$. Mostre que $X$ está quase
certamente concentrada em uma progressão aritmética $a + \frac{2\pi}{\xi_0}\Z$ \emph{(escreva
$\varphi_X(\xi_0) = \eu^{\iu\theta}$ e calcule $\E[1 -
\cos(\xi_0X - \theta)]$)}. Deduza que, se $X$ tem \omterm{ex:b3:lebesgue:gamma}{densidade}, então
$\abs{\varphi_X(\xi)} < 1$ para todo $\xi \neq 0$.
\end{exercise}

\begin{exercise}[$\star\star$]\label{exo:b3:clt:3}
Sejam $X \sim \mathcal N(m_1, \sigma_1^2)$ e $Y \sim
\mathcal N(m_2, \sigma_2^2)$ \omterm{def:b3:probability:independence}{independentes}. Mostre que $X + Y
\sim \mathcal N(m_1 + m_2, \sigma_1^2 + \sigma_2^2)$ e, mais geralmente,
que a família \omterm{def:b3:clt:gaussianvector}{gaussiana} é estável por somas \omterm{def:b3:probability:independence}{independentes} e por
aplicações afins. Em contraste: a soma de duas \omterm{def:b3:clt:gaussianvector}{gaussianas}
\emph{dependentes} é sempre \omterm{def:b3:clt:gaussianvector}{gaussiana}? (\cref{exo:b3:clt:9}.)
\end{exercise}

\begin{exercise}[$\star\star$]\label{exo:b3:clt:4}
(a) Demonstre a equivalência da~\cref{def:b3:clt:cid}: se $\E f(X_n) \to \E f(X)$ para toda
$f$ \omterm{def:b3:topology:continuity}{contínua} limitada, então $F_{X_n}(t) \to F_X(t)$ nos pontos de \omterm{def:b3:topology:continuity}{continuidade}
\emph{(comprima $\mathbf 1_{\intoc{-\infty}t}$ entre duas rampas \omterm{def:b3:topology:continuity}{contínuas} em escada)}; e
reciprocamente \emph{(aproxime uma $f$ \omterm{def:b3:topology:continuity}{contínua} limitada por somas de
funções-rampa, ou condicione em uma grade fina de pontos de
\omterm{def:b3:topology:continuity}{continuidade})} --- a recíproca pode ser tratada primeiro para $f$
uniformemente \omterm{def:b3:topology:continuity}{contínua} e, depois, em geral.
(b) Mostre que $X_n \Rightarrow c$ (uma constante) implica $X_n
\to c$ em probabilidade.
\end{exercise}

\begin{exercise}[$\star\star$]\label{exo:b3:clt:5}
(Lei dos eventos raros) Sejam $X_n \sim \mathcal B(n, p_n)$ com $np_n \to \lambda > 0$. Mostre, via \omterm{def:b3:clt:cf}{funções
características} e o~\cref{thm:b3:clt:levy}, que $X_n \Rightarrow \mathcal
P(\lambda)$. Verificação numérica de
sanidade: compare $\P(X = 0)$ para $\mathcal B(100, 0.02)$ e $\mathcal P(2)$.
\end{exercise}

\begin{exercise}[$\star\star$]\label{exo:b3:clt:6}
(a) Um dado honesto é lançado $n = 1000$ vezes; aproxime a probabilidade de o
total ultrapassar $3600$ (média $3500$, variância por lançamento
$\frac{35}{12}$).
(b) Para $S \sim \mathcal B(100, \frac12)$, aproxime $\P(45 \leq S \leq 55)$ pelo teorema central do limite com a
correção de \omterm{def:b3:topology:continuity}{continuidade} ($\pm\frac12$) e comente o efeito da correção.
\end{exercise}

\begin{exercise}[$\star\star$]\label{exo:b3:clt:7}
(Erro de \omterm{ex:b3:probability:sllnapps}{Monte Carlo}) No cenário do~\cref{pb:b3:probability:1}, questão 11, com
$g \in
L^2(\intcc01^d)$, sejam $\sigma^2 = \V(g(U_1))$ e $I = \int
g$. Mostre que
\[
\sqrt n\,\Bigl(\frac1n\sum_{k\leq n}g(U_k) - I\Bigr)
\Longrightarrow \mathcal N(0, \sigma^2),
\]
e deduza a barra de erro assintótica de $95\%$, $\pm
1.96\,\sigma/\sqrt n$ ---
\omterm{def:b3:probability:independence}{independente} da dimensão $d$. Compare com a regra determinística do
ponto médio em dimensão $d$ (erro $\sim n^{-2/d}$ para integrandos $\mathcal C^2$): a
partir de que dimensão a amostragem aleatória vence?
\end{exercise}

\begin{exercise}[$\star\star\star$]\label{exo:b3:clt:8}
(Slutsky) Suponha $X_n \Rightarrow X$ e $Y_n \to c$ em probabilidade ($c$ constante).
Mostre que $X_n + Y_n \Rightarrow X +
c$ e $Y_nX_n \Rightarrow cX$. \emph{(Trabalhe com \omterm{def:b3:clt:cf}{funções características} e
com a cota $\abs{\E\eu^{\iu\xi
(X_n+Y_n)} - \eu^{\iu\xi c}\E\eu^{\iu\xi X_n}} \leq
\E\abs{\eu^{\iu\xi(Y_n - c)} - 1}$, separando em $\abs{Y_n - c}
\leq \delta$.)} Aplicação: no~\cref{ex:b3:clt:confidence},
justifique a substituição do $\sigma = \sqrt{p(1-p)}$ desconhecido por $\sqrt{\hat p_n(1 - \hat
p_n)}$.
\end{exercise}

\begin{exercise}[$\star\star\star$]\label{exo:b3:clt:9}
Sejam $X \sim \mathcal N(0,1)$ e $\varepsilon$ \omterm{def:b3:probability:independence}{independentes} com $\P(\varepsilon = \pm1) = \frac12$; ponha $Y =
\varepsilon X$.
(a) Mostre que $Y \sim \mathcal N(0,1)$ e $\operatorname{Cov}(X, Y) = 0$.
(b) Mostre que $X$ e $Y$ \emph{não} são \omterm{def:b3:probability:independence}{independentes} e que
$(X, Y)$ não é um \omterm{def:b3:clt:gaussianvector}{vetor gaussiano} \emph{(calcule $\P(X + Y =
0)$)}.
(c) Moral: o~\cref{thm:b3:clt:gaussianvector}(2) exige gaussianidade conjunta ---
``\omterm{def:b3:clt:gaussianvector}{gaussianas} não correlacionadas'', por si só, nada demonstram.
\end{exercise}

\begin{exercise}[$\star\star$]\label{exo:b3:clt:10}
A lei de Cauchy tem \omterm{ex:b3:lebesgue:gamma}{densidade} $\frac1{\pi(1 + x^2)}$.
(a) Mostre que sua \omterm{def:b3:clt:cf}{função característica} é $\eu^{-\abs\xi}$ (\cref{exo:b3:fouriertransform:1} e
inversão).
(b) Mostre que, se $X_1, \dots, X_n$ são i.i.d.\ de Cauchy, então $\frac{S_n}n$ é de novo de
Cauchy --- a \emph{mesma} lei: a média jamais se concentra.
(c) Reconcilie com as leis dos grandes números e com o teorema central
do limite: que hipóteses falham? (Calcule $\E\abs{X_1}$.)
\end{exercise}

\begin{exercise}[$\star\star$]\label{exo:b3:clt:11}
(Leis estáveis em embrião) Sejam $(X_n)$ i.i.d.\ de Cauchy padrão
(\cref{exo:b3:clt:10}).
(a) Mostre que, para qualquer $a, b > 0$, $aX_1 + bX_2$ tem a lei de $(a + b)X_1$: a família
de Cauchy é \emph{estritamente estável} de índice $1$.
(b) Mostre que a família \omterm{def:b3:clt:gaussianvector}{gaussiana} é estritamente estável de índice
$2$: $aX_1 + bX_2 \sim \sqrt{a^2 + b^2}\,X_1$ para $X_i$ i.i.d.\ $\mathcal N(0,1)$.
(c) Explique, via \omterm{def:b3:clt:cf}{funções características} da forma $\eu^{-c\abs\xi^\alpha}$, por que a
estabilidade de índice $\alpha$ força a normalização $n^{1/\alpha}$ para as
somas, e o que isso diz sobre as bacias de atração do teorema central do
limite: que somas i.i.d.\ podem convergir, após normalização afim, a uma
lei de Cauchy em vez de a uma \omterm{def:b3:clt:gaussianvector}{gaussiana}?
\end{exercise}

\begin{exercise}[$\star\star$]\label{exo:b3:clt:12}
(A função de distribuição empírica) Sejam $(X_n)$ i.i.d.\ com função
de distribuição $F$, e $F_n(t) =
\frac1n\#\{k \leq n : X_k \leq t\}$.
(a) Fixe $t$. Mostre que $n F_n(t) \sim \mathcal B(n, F(t))$, que $F_n(t) \to F(t)$ q.c.\ (\cref{thm:b3:probability:slln}) e
que
\[
\sqrt n\,\bigl(F_n(t) - F(t)\bigr) \Longrightarrow
\mathcal N\bigl(0,\ F(t)(1 - F(t))\bigr) .
\]
(b) Em que $t$ a variância assintótica é máxima? Interprete: a
mediana é o ponto em que uma distribuição empírica é mais difícil de
fixar.
(c) Para $F$ \omterm{def:b3:topology:continuity}{contínua}, mostre que a lei de $\sup_t\abs{F_n(t) - F(t)}$ não depende de
$F$ \emph{(reduza a variáveis uniformes via o~\cref{exo:b3:probability:1})} --- o
milagre livre de distribuição por trás do teste de
Kolmogorov--Smirnov; não se pede calcular essa lei.
\end{exercise}

\section{Problema: a demonstração de Lindeberg do teorema central do
limite, com taxa}

\begin{problem}[{Problema de fim de semana --- o método da
substituição}]
\label{pb:b3:clt:1}
Lindeberg (1922) demonstrou o teorema central do limite por uma ideia de
desarmante simplicidade: \emph{trocar as parcelas, uma de cada vez, por
\omterm{def:b3:clt:gaussianvector}{gaussianas}} e controlar cada troca por uma expansão de Taylor. O método
não precisa de análise de Fourier, produz uma taxa de erro explícita e,
hoje, move demonstrações de universalidade por toda a teoria de
probabilidade. Sejam $(X_i)$ i.i.d., centradas, $\V(X_1) = 1$, com $\beta = \E\abs{X_1}^3 < \infty$;
sejam $(N_i)$ i.i.d.\ $\mathcal N(0,1)$, \omterm{def:b3:probability:independence}{independentes} das $X_i$
(existência: \cref{thm:b3:probability:existence}). Ponha
\[
T_n = \frac{X_1 + \dots + X_n}{\sqrt n},
\qquad
G_n = \frac{N_1 + \dots + N_n}{\sqrt n} \sim \mathcal N(0,1).
\]

\medskip
\noindent\textbf{Parte I --- A identidade de substituição.} Fixe $f
\in \mathcal C^3_b(\R)$
(três derivadas \omterm{def:b3:topology:continuity}{contínuas} e limitadas; $M_3 = \sup\abs{f'''}$). Para $0 \leq i \leq n$, defina as
somas híbridas
\[
H_i = \frac{X_1 + \dots + X_i + N_{i+1} + \dots +
N_n}{\sqrt n},
\]
de modo que $H_n = T_n$ e $H_0 = G_n$.
\begin{enumerate}
  \item Escreva $H_i = W_i + \frac{X_i}{\sqrt n}$ e $H_{i-1}
        = W_i + \frac{N_i}{\sqrt n}$ com $W_i =
        \frac{1}{\sqrt n}\bigl(\sum_{j<i}X_j +
        \sum_{j>i}N_j\bigr)$, e note que $W_i$ é
        \omterm{def:b3:probability:independence}{independente} do par $(X_i, N_i)$. Justifique.
  \item Taylor com resto integral ou de Lagrange: para reais $w, h$
        quaisquer:
        \[
        \Bigl|f(w + h) - f(w) - f'(w)h -
        \tfrac12f''(w)h^2\Bigr| \leq
        \frac{M_3\,\abs h^3}{6} .
        \]
  \item Aplique a questão 2 duas vezes ($h = \frac{X_i}{\sqrt n}$ e $h = \frac{N_i}{\sqrt n}$ em $w = W_i$), tome
        \omterm{def:b3:probability:space}{esperanças} e use a \omterm{def:b3:probability:independence}{independência} mais a coincidência dos dois
        primeiros momentos de $X_i$ e $N_i$ para mostrar que
        \[
        \bigl|\E f(H_i) - \E f(H_{i-1})\bigr|
        \leq \frac{M_3}{6}\cdot
        \frac{\beta + \gamma}{n^{3/2}},
        \qquad \gamma = \E\abs{N_1}^3 =
        \frac{2\sqrt2}{\sqrt\pi} .
        \]
  \item Telescope em $i$ e conclua a \emph{cota de Lindeberg}:
        \[
        \bigl|\E f(T_n) - \E f(G_n)\bigr| \leq
        \frac{M_3\,(\beta + \gamma)}{6\,\sqrt n} .
        \]
\end{enumerate}

\medskip
\noindent\textbf{Parte II --- De $f$ suave ao teorema central do
limite.}
\begin{enumerate}[resume]
  \item Mostre que $\E f(T_n) \to \E f(N)$ para toda $f \in
        \mathcal C_b^3$, e eleve a toda $f$ \omterm{def:b3:topology:continuity}{contínua}
        limitada: dada uma tal $f$ e $\varepsilon$, construa $f_\varepsilon \in \mathcal C^3_b$ com $\norm{f - f_\varepsilon}_\infty \leq \varepsilon$
        em um intervalo grande --- por exemplo, convolua $f$ com uma
        \omterm{def:b3:lp:mollifier}{função de corte} $\mathcal C^\infty$ (\cref{thm:b3:lp:regularization}) --- e trate as caudas por
        tensão ($\V(T_n) = 1$ e Chebyshev). Conclua $T_n \Rightarrow \mathcal N(0, 1)$: o teorema central do
        limite, redemonstrado.
  \item Onde a demonstração usou que as $X_i$ são
        \emph{identicamente} distribuídas? Mostre que quase não usou:
        enuncie e demonstre a versão para $X_i$ \omterm{def:b3:probability:independence}{independentes},
        centradas e não idênticas, com $\sum_i\V(X_i) = s_n^2$ e terceiros momentos,
        obtendo o erro $\frac{M_3}{6s_n^3}\sum_i\bigl(\E\abs{X_i}^3
        + \V(X_i)^{3/2}\gamma\bigr)$ --- o verdadeiro teorema de Lindeberg em sua
        forma de Lyapunov.
\end{enumerate}

\medskip
\noindent\textbf{Parte III --- Dividendos quantitativos.}
\begin{enumerate}[resume]
  \item (Funções de distribuição) Seja $t \in \R$ e aproxime $\mathbf 1_{\intoc{-\infty}t}$ por cima e
        por baixo por rampas $\mathcal C^3_b$ de largura $\delta$ (construa-as,
        com $M_3 = O(\delta^{-3})$). Combinando com a Parte I, deduza a cota de dois
        termos
        \[
        \sup_{t\in\R}\,\bigl|\P(T_n \leq t) -
        \Phi(t)\bigr| \;\leq\;
        \frac{C_1(\beta + \gamma)}{\delta^3\sqrt n} +
        C_2\,\delta
        \qquad (\text{todo } \delta > 0),
        \]
        com constantes explícitas (o termo $C_2\delta$ usa que $\Phi$ tem
        \omterm{ex:b3:lebesgue:gamma}{densidade} limitada por $\frac1{\sqrt{2\pi}}$) e otimize $\delta \sim
        n^{-1/8}$ para obter uma taxa
        uniforme de ordem $n^{-1/8}$. (O $n^{-1/2}$ ótimo --- Berry--Esseen ---
        exige ferramentas mais finas; o que importa aqui é obter uma
        taxa \emph{explícita} por substituição elementar.)
  \item (De Moivre--Laplace, quantificado) Especialize a $X_i = 2B_i - 1$ (sinais
        de moedas honestas): compare a conclusão com a estimativa local
        do~\cref{pb:b3:product:1}, questão 7 --- o que cada método dá e o outro
        não dá?
  \item (Universalidade) Explique em um parágrafo por que o método da
        substituição mostra mais do que o teorema central do limite:
        qualquer estatística da forma $\E f(\text{soma})$ com $f$ suave é
        insensível, na ordem $n^{-1/2}$, à \emph{lei inteira} das parcelas
        para além de seus dois primeiros momentos --- o ``princípio de
        invariância'' que sustenta os resultados modernos de
        universalidade (matrizes aleatórias, polinômios aleatórios), do
        qual o teorema central do limite é a primeira instância.
\end{enumerate}

\medskip
\noindent\textbf{Parte IV --- Suavização, levada adiante: taxas
melhores.} A perda de $n^{-1/2}$ (com $f$ suave) para $n^{-1/8}$ (funções de
distribuição) veio de cobrar $f'''$ na norma do supremo. Os híbridos
podem reparar parte disso: eles contêm parcelas \omterm{def:b3:clt:gaussianvector}{gaussianas}, e as
\omterm{def:b3:clt:gaussianvector}{gaussianas} \emph{suavizam}.
\begin{enumerate}[resume]
  \item (Uma \omterm{def:b3:clt:gaussianvector}{gaussiana} escondida) Para $1 \leq i \leq n - 1$, $h =
        \frac{X_i}{\sqrt n}$ ou $\frac{N_i}{\sqrt n}$, e $\theta \in \intcc01$,
        escreva $W_i + \theta h = A +
        Z$ com $Z = \frac{N_{i+1} + \dots + N_n}{\sqrt
        n}$. Mostre que $Z \sim \mathcal N\bigl(0,
        \frac{n-i}n\bigr)$ é \omterm{def:b3:probability:independence}{independente} do par
        $(A,
        h)$ e deduza, para toda $g \in
        L^1(\R)$ \omterm{def:b3:topology:continuity}{contínua},
        \[
        \E\bigl[\abs h^3\,\abs{g(W_i + \theta h)}\bigr]
        \;\leq\; \sqrt{\frac{n}{2\pi(n - i)}}\;
        \norm{g}_{L^1}\;\E\abs h^3 .
        \]
  \item Combine a questão 10 com a forma integral do resto de Taylor,
        \[
        f(w + h) = f(w) + f'(w)h + \tfrac12f''(w)h^2 +
        \int_0^1\frac{(1 - \theta)^2}2\,f'''(w + \theta
        h)\,h^3\,\dd\theta,
        \]
        para refazer as questões 3--4: para $f \in \mathcal C^3_b$ com, além disso,
        $f''' \in L^1(\R)$,
        \[
        \bigl|\E f(T_n) - \E f(G_n)\bigr| \leq
        \frac{\beta + \gamma}{3\sqrt{2\pi}}\cdot
        \frac{\norm{f'''}_{L^1}}{\sqrt n} +
        \frac{M_3(\beta + \gamma)}{6\,n^{3/2}}
        \]
        \emph{(a questão 10 cuida das trocas $i \leq n - 1$ --- use $\sum_{m=1}^{n-1}m^{-1/2} \leq 2\sqrt n$ --- e a
        cota grosseira da questão 3 cuida da última)}. Verifique que as
        rampas da questão 7 satisfazem $\norm{\psi_\delta'''}_{L^1} = K_1\delta^{-2}$ enquanto $M_3 = K\delta^{-3}$, insira-as e
        otimize $\delta$: a taxa uniforme para funções de
        distribuição melhora para $O(n^{-1/6})$.
  \item (Coincidindo mais um momento) Suponha, além disso, $\E
        X_1^3 = 0$ e
        $\beta_4 = \E X_1^4 < \infty$. Calcule $\E N_1^3$ e $\E N_1^4$, expanda até a quarta ordem e
        demonstre pelas mesmas linhas que a taxa para funções de
        distribuição se torna $O(n^{-1/4})$ \emph{(agora $\norm{\psi_\delta^{(4)}}_{L^1} =
        K_2\delta^{-3}$ e $M_4 = K'\delta^{-4}$; escolha
        $\delta = n^{-1/4}$)}.
  \item (A obstrução) Suponha que os $k$ primeiros momentos de
        $X_1$ coincidam com os \omterm{def:b3:clt:gaussianvector}{gaussianos} ($k = 2$ sempre;
        $k = 3$ exatamente quando $\E X_1^3 = 0$; $k \geq 4$ essencialmente nunca,
        pois $\E N_1^4 = 3$). Verifique que o esquema das questões 10--12 entrega
        a taxa $n^{-(k-1)/(2k+2)}$ para funções de distribuição, equilibrando $\delta^{-k}n^{-(k-1)/2}$
        contra $\delta$, e observe que o expoente só se aproxima do
        valor $\frac12$ de Berry--Esseen quando $k \to
        \infty$. Explique em algumas
        frases por que o método da substituição satura: cada troca é
        cobrada em valor absoluto, ao passo que a rota de Fourier
        (desigualdade de suavização de Esseen) explora a oscilação da
        diferença de \omterm{def:b3:clt:cf}{funções características} e atinge $C\beta n^{-1/2}$ apenas com
        três momentos.
\end{enumerate}

\medskip
\noindent\textbf{Parte V --- Duas dimensões: o teorema central do
limite multidimensional, por substituição.} Sejam agora as $X_i$
\emph{vetores} aleatórios i.i.d.\ centrados de $\R^2$, de matriz de
covariância $\Sigma$ e $\beta' = \E\norm{X_1}^3 <
\infty$ (norma euclidiana).
\begin{enumerate}[resume]
  \item (\omterm{def:b3:clt:gaussianvector}{Vetores gaussianos}, na ordem certa) Diagonalize $\Sigma =
        PDP^{\mathsf T}$
        (\cref{exo:b3:submanifolds:8}) e ponha $C = P\sqrt DP^{\mathsf T}$. Para $Z = (Z^1,
        Z^2)$ um par de \omterm{def:b3:clt:gaussianvector}{gaussianas}
        padrão \omterm{def:b3:probability:independence}{independentes} (\cref{thm:b3:probability:existence}), mostre que $N
        = CZ$ é
        um \omterm{def:b3:clt:gaussianvector}{vetor gaussiano} (\cref{def:b3:clt:gaussianvector}) de média $0$, covariância
        $\Sigma$ e $\gamma' = \E\norm N^3 <
        \infty$; e que $G_n = \frac{N_1 + \dots +
        N_n}{\sqrt n}$ tem lei $\mathcal N(0, \Sigma)$ \emph{exatamente},
        para cópias i.i.d.\ $N_i$.
  \item (Taylor em duas variáveis) Para $f \colon \R^2 \to
        \R$ de classe $\mathcal C^3$ com $M_3 =
        \max_{\abs\alpha = 3}\sup\abs{\partial^\alpha f} <
        \infty$,
        demonstre
        \[
        \Bigl|f(w + h) - f(w) - \langle\nabla f(w),
        h\rangle - \tfrac12\langle h, D^2f(w)\,h\rangle
        \Bigr| \leq \frac{M_3}6\,\bigl(\abs{h_1} +
        \abs{h_2}\bigr)^3 \leq \frac{\sqrt2\,M_3}3\,
        \norm h^3
        \]
        \emph{(estude $t \mapsto f(w + th)$ em $\intcc01$)}.
  \item (O teorema central do limite em $\R^2$) Rode o esquema de
        substituição nos híbridos vetoriais $H_i$: mostre que os
        termos de primeira e de segunda ordem se cancelam (médias e
        covariâncias coincidem), telescope e eleve como na questão 5
        (tensão a partir de $\E\norm{T_n}^2 =
        \operatorname{tr}\Sigma$; regularização agora em $\R^2$,
        \cref{thm:b3:lp:regularization}) para concluir: para toda $f \colon
        \R^2 \to \R$ \omterm{def:b3:topology:continuity}{contínua} limitada,
        \[
        \E\,f\Bigl(\frac{X_1 + \dots + X_n}{\sqrt n}\Bigr)
        \longrightarrow \E\,f(N), \qquad N \sim \mathcal
        N(0, \Sigma) :
        \]
        o~\cref{thm:b3:clt:multiclt} em dimensão $2$, com taxa para $f$ suave
        e sem análise de Fourier.
  \item (Cramér--Wold e uma flutuação conjunta) Deduza que $\langle t, \frac{S_n}{\sqrt n}\rangle
        \Rightarrow \mathcal N(0, t^{\mathsf T}\Sigma t)$ para
        todo $t \in \R^2$ fixado. Aplicação: para $(\xi_i)$ reais i.i.d., centradas,
        com $\E\xi_1^2 = 1$, $\E\xi_1^6 < \infty$ (de modo que a Parte V se aplica a $V_i = (\xi_i, \xi_i^2 - 1)$), mostre
        que
        \[
        \frac1{\sqrt n}\Bigl(\sum_{i\leq n}\xi_i,\
        \sum_{i\leq n}(\xi_i^2 - 1)\Bigr) \Longrightarrow
        \mathcal N\Bigl(0, \begin{pmatrix} 1 & \E\xi_1^3\\
        \E\xi_1^3 & \E\xi_1^4 - 1\end{pmatrix}\Bigr) :
        \]
        a média empírica e o segundo momento empírico flutuam
        conjuntamente de maneira \omterm{def:b3:clt:gaussianvector}{gaussiana} --- independentemente, no
        limite, se, e somente se, $\E\xi_1^3 = 0$ (\cref{thm:b3:clt:gaussianvector}).
\end{enumerate}

\medskip
\noindent\textbf{Parte VI --- O método delta.}
\begin{enumerate}[resume]
  \item Sejam $(\hat\theta_n)$ \omterm{def:b3:probability:space}{variáveis aleatórias} com $\sqrt n(\hat\theta_n - \theta) \Rightarrow
        \mathcal N(0, \sigma^2)$ para um parâmetro
        real $\theta$, e seja $g$ derivável em $\theta$.
        Demonstre o \emph{método delta}:
        \[
        \sqrt n\bigl(g(\hat\theta_n) - g(\theta)\bigr)
        \Longrightarrow \mathcal N\bigl(0,
        g'(\theta)^2\sigma^2\bigr)
        \]
        \emph{(escreva $g(x) - g(\theta) = (g'(\theta) +
        \eta(x))(x - \theta)$ com $\eta \to 0$ em $\theta$; mostre que $\hat\theta_n \to \theta$
        e, depois, $\eta(\hat\theta_n) \to 0$, em probabilidade; termine com Slutsky,
        \cref{exo:b3:clt:8}, e com o~\cref{exo:b3:clt:4}(b))}.
  \item Aplicações. (a) Para $(\xi_i)$ reais i.i.d.\ de média $\mu$ e
        variância $\sigma^2$, e $\bar X_n =
        \frac1n\sum_{i\leq n}\xi_i$: mostre $\sqrt n(\bar
        X_n^2 - \mu^2) \Rightarrow \mathcal N(0,
        4\mu^2\sigma^2)$ quando $\mu \neq 0$, e que, para
        $\mu = 0$, o enunciado correto vive em outra escala: $n\bar X_n^2 \Rightarrow \sigma^2N^2$ com $N \sim \mathcal N(0,1)$
        (identifique a função de distribuição do limite). (b)
        (Estabilização da variância) Para $\hat p_n$ a frequência de sucessos
        de uma amostra $\mathcal B(1, p)$, $p \in \intoo01$: mostre que $g(p) = \arcsin\sqrt p$ satisfaz
        \[
        \sqrt n\,\bigl(g(\hat p_n) - g(p)\bigr)
        \Longrightarrow \mathcal N\Bigl(0, \frac14\Bigr)
        \]
        \emph{qualquer} que seja $p$ --- uma barra de erro
        assintótica livre do parâmetro desconhecido; compare com
        o~\cref{ex:b3:clt:confidence}.
\end{enumerate}

\medskip
\noindent\textbf{Parte VII --- Poisson, pelo mesmo método: o teorema de
Le Cam.} A substituição conhece uma segunda classe de universalidade:
somas de muitos eventos \omterm{def:b3:probability:independence}{independentes} \emph{raros}. Para leis em
$\N$, a distância certa é a \emph{variação total},
\[
d_{\mathrm{TV}}(\mu, \nu) = \sup_{A\subseteq\N}\,
\abs{\mu(A) - \nu(A)} .
\]
\begin{enumerate}[resume]
  \item Mostre que $d_{\mathrm{TV}}(\mu, \nu) =
        \frac12\sum_{k\geq0}\abs{\mu(\{k\}) -
        \nu(\{k\})}$ e demonstre a cota de acoplamento: para
        \emph{qualquer} par $(X, Y)$ de \omterm{def:b3:probability:space}{variáveis aleatórias} de leis
        $\mu$ e $\nu$ no mesmo espaço, $d_{\mathrm{TV}}(\mu, \nu) \leq \P(X \neq Y)$.
  \item Calcule exatamente, para $p \in \intoo01$:
        \[
        d_{\mathrm{TV}}\bigl(\mathcal B(1, p), \mathcal
        P(p)\bigr) = p\bigl(1 - \eu^{-p}\bigr) \leq p^2 .
        \]
  \item (Le Cam, por substituição) Sejam $X_i \sim \mathcal B(1,
        p_i)$ e $Y_i \sim \mathcal P(p_i)$, com as $2n$
        variáveis \omterm{def:b3:probability:independence}{independentes}; $S = X_1 + \dots + X_n$, e recorde $Y_1 + \dots + Y_n \sim \mathcal P(\lambda)$ com $\lambda = \sum_ip_i$
        (\cref{exo:b3:clt:1}). Troque uma coordenada de cada vez nos híbridos
        inteiros $H_i = Y_1 + \dots + Y_i + X_{i+1} + \dots
        + X_n$: mostre, para todo $A \subseteq \N$,
        \[
        \abs{\P(H_{i-1} \in A) - \P(H_i \in A)} \leq
        d_{\mathrm{TV}}\bigl(\mathcal B(1, p_i), \mathcal
        P(p_i)\bigr),
        \]
        e conclua a \emph{desigualdade de Le Cam}:
        \[
        d_{\mathrm{TV}}\bigl(\text{lei de } S,\ \mathcal
        P(\lambda)\bigr) \leq \sum_{i=1}^np_i^2 .
        \]
  \item Dividendos. (a) Para $p_i = \frac\lambda n$: a cota vale $\frac{\lambda^2}n$ --- a lei dos
        eventos raros (\cref{exo:b3:clt:5}) elevada a uma taxa explícita,
        uniforme sobre todos os eventos e válida também para $p_i$
        desiguais. (b) $500$ cartas são entregues, cada uma
        extraviando-se de maneira \omterm{def:b3:probability:independence}{independente} com probabilidade $\frac1{500}$:
        limite o erro do modelo de Poisson de parâmetro $1$ e estime
        a probabilidade de nenhuma carta se extraviar. (c) Encerre o
        problema: compare as duas classes de universalidade aqui
        encontradas --- a \omterm{def:b3:clt:gaussianvector}{gaussiana} (muitas contribuições pequenas e
        espalhadas; dois momentos coincidentes; Taylor) e a de Poisson
        (muitas contribuições raras; uma média coincidente; um
        acoplamento exato em variação total) --- e o único método de
        substituição por trás de ambas.

  \item (Erro relativo e a transformação logarítmica) Sejam $(X_n)$
        i.i.d., positivas, de média $\mu > 0$, variância $\sigma^2$, e seja $\bar X_n$
        a média empírica. Mostre pelo método delta que
        \[
        \sqrt n\,\bigl(\ln\bar X_n - \ln\mu\bigr)
        \Longrightarrow
        \mathcal N\Bigl(0,\ \frac{\sigma^2}{\mu^2}\Bigr) :
        \]
        o parâmetro assintótico de $\ln\bar X_n$ é o \emph{coeficiente de
        variação} $\sigma/\mu$ --- um erro relativo, livre de escala. Deduza um
        intervalo de confiança de $95\%$ para $\mu$ na forma
        multiplicativa $\bar X_n\cdot\eu^{\pm1.96\,\sigma/(\mu\sqrt
        n)}$ e explique quando ele é preferível ao
        aditivo.
  \item (O terceiro momento comanda o erro) Para $X \sim$
        Bernoulli($p$) centrada, calcule $\E\bigl[(X -
        p)^3\bigr] = p(1-p)(1-2p)$. Usando a análise da
        Parte IV (o erro de substituição é comandado pelos terceiros
        momentos), explique por que a aproximação normal de $\mathcal B(n, p)$ é
        assimétrica para $p \neq \frac12$ --- excedendo de um lado e ficando aquém
        do outro --- e por que $p = \frac12$ desfruta da taxa mais rápida com
        momentos coincidentes. Verifique numericamente o sinal da
        assimetria em $\mathcal
        B(20, 0.1)$ contra $\mathcal N(2, 1.8)$: compare $\P(S = 0) = 0.9^{20}$ com a massa
        \omterm{def:b3:clt:gaussianvector}{gaussiana} de $\intoo{-\infty}{0.5}$.
\end{enumerate}
\end{problem}
